\chapter{第二部分：Spark 大数据分析}

本部分在第一部分的 CCE 集群上，通过 Spark Operator 提交 PySpark 作业，对豆瓣电影
数据集（OBS 存储，s3a 接口读取，约 6.7 万条记录）完成环境部署、数据清洗、
SQL 统计分析与性能对比。

\section{A-0：环境部署}

\subsection{Spark Operator 安装}

使用离线 Helm Chart 安装 Spark Operator，并通过 \lstinline!--set! 将控制器/Webhook
镜像指向 SWR group17（避免国内无法访问 ghcr.io）：

\begin{lstlisting}[language=bash,caption={安装 Spark Operator},captionpos=b]
helm install spark-op ./spark-operator -n spark-operator \
  --set image.registry=swr.cn-north-4.myhuaweicloud.com \
  --set image.repository=group17/spark-operator --set image.tag=2.5.0 \
  --set image.pullSecrets[0].name=default-secret \
  --set spark.jobNamespaces={default}
\end{lstlisting}

\subsection{关键问题：Executor 因 CPU 请求不足无法调度}

首次提交作业时 Driver 持续 Pending，\lstinline!describe! 显示
\lstinline!Insufficient cpu!。原因是 2 个 2\,vCPU 节点被 CCE 系统组件占用了约
90\% 的 CPU \textbf{请求}额度（实际使用率仅个位数），而 Driver 默认按
\lstinline!cores: 1! 申请 1000m CPU。解决办法是用 \lstinline!coreRequest! 把
Kubernetes 的 CPU 请求与 Spark 逻辑核心解耦：Driver \lstinline!coreRequest: "100m"!、
Executor \lstinline!coreRequest: "50m"!，并临时把 Web 应用缩容到 0 释放额度。
调整后 Driver 与 2 个 Executor 均成功 Running。

\subsection{WordCount 验证}

提交 \lstinline!wordcount.py!（读取 \lstinline!s3a://group17/data/sample.txt!），
Driver 状态变为 Completed，输出 Top10 词频，验证 Spark Operator + OBS s3a 读取链路
正常：

\begin{lstlisting}[caption={WordCount 输出},captionpos=b]
Top 10 words: [('spark', 300), ('hello', 200), ('data', 150),
('cloud', 100), ('is', 50), ('design', 50), ('world', 50), ...]
\end{lstlisting}

提交豆瓣分析作业后，\lstinline!kubectl get pods -n default! 可见 1 个 Driver 与
2 个 Executor Pod（图~\ref{fig:sparkpods}）。

\shot{spark_driver_executor.png}{Spark Driver + 2 个 Executor Pod 均 Running}{fig:sparkpods}

\section{A-1：数据清洗}

用 PySpark 读取 OBS 上的 CSV（开启 \lstinline!multiLine! 以正确解析含换行的 summary
字段），强制转换 year/rating\_score/rating\_count 等字段类型，打印 Schema 与前 5 行，
原始记录 \textbf{67{,}132} 条。各字段缺失值比例如表~\ref{tab:miss}（注意
rating\_score 为 0 表示"评分人数不足、豆瓣未出分"，等价于缺失，缺失率高达 65\%）。

\begin{table}[htbp]\centering
\caption{各字段缺失值比例}\label{tab:miss}
\begin{tabular}{lc lc}
\toprule
字段 & 缺失比例 & 字段 & 缺失比例 \\
\midrule
year         & 3.00\%  & directors    & 13.15\% \\
rating\_score & 65.06\% & summary      & 28.56\% \\
genres       & 8.02\%  & 其余字段     & 0.00\% \\
\bottomrule
\end{tabular}
\end{table}

\textbf{两种不同的缺失值处理策略：}
\begin{itemize}
  \item \textbf{year（时间维度关键字段）—— dropna 删行}：年份缺失或越界
  （非 1920--2025）的记录无法用于时间趋势分析，直接删除。共删除 2{,}013 行，
  清洗后剩 \textbf{65{,}119} 行。
  \item \textbf{rating\_score（分析核心指标）—— fillna 均值填充}：直接删除 65\%
  的未评分电影会损失过多样本，故新增 \lstinline!rating_filled! 列，用"已评分电影
  均分 6.55"填充缺失/为 0 的值，保留样本量用于计数类统计；而在评分高低相关查询中
  再过滤 \lstinline!rating_score > 0!（共 22{,}724 部已评分电影）以保证均值不被
  填充值干扰。
\end{itemize}

数值字段基本统计（describe）显示：year 均值 1998、rating\_score（含 0）均值仅 2.41
而 rating\_filled 均值 6.55、rating\_count 最大约 199 万（《肖申克的救赎》）。
清洗前后行数对比见图~\ref{fig:sparka1}。

\shot{spark_a1_clean.png}{A-1 数据清洗：Schema、缺失值比例、清洗前后行数与统计}{fig:sparka1}

\section{A-2：Spark SQL 统计分析}

基于已评分电影视图，用 Spark SQL 完成 5 个查询，覆盖 GROUP BY 聚合、ORDER BY Top-N、
时间维度趋势、窗口函数与 JOIN。

\subsection{Q1：各电影类型数量与平均评分（GROUP BY）}

\begin{table}[htbp]\centering
\caption{各类型电影数量与平均评分 Top（部分）}\label{tab:q1}
\begin{tabular}{lcc lcc}
\toprule
类型 & 数量 & 均分 & 类型 & 数量 & 均分 \\
\midrule
剧情 & 11662 & 6.83 & 动画 & 1531 & 7.25 \\
喜剧 & 5079  & 6.61 & 家庭 & 705  & 7.22 \\
爱情 & 3610  & 6.50 & 传记 & 567  & 7.27 \\
动作 & 3132  & 6.13 & 恐怖 & 2074 & 5.48 \\
\bottomrule
\end{tabular}
\end{table}

\textbf{分析}：剧情片数量最多（1.17 万部）且均分较高（6.83），是豆瓣电影的绝对主力；
传记（7.27）、动画（7.25）、家庭（7.22）类均分最高，说明这类影片质量普遍较好、口碑
稳定；而恐怖（5.48）、惊悚（5.90）类均分明显偏低，反映此类商业片良莠不齐、整体口碑
一般。类型与评分的关系为内容生产提供了参考。

\subsection{Q2：高分电影 Top10（ORDER BY Top-N）}

筛选评分人数 $\geq 5$ 万的电影按评分降序，Top10 为：肖申克的救赎(9.7)、霸王别姬(9.6)、
控方证人(9.6)、阿甘正传(9.5)、美丽人生(9.5)、辛德勒的名单(9.5)、人生果实(9.5)、
这个杀手不太冷(9.4)、千与千寻(9.4)、泰坦尼克号(9.4)。

\textbf{分析}：加评分人数门槛可有效过滤"小众高分"噪声，得到的高分榜单与公认经典高度
一致。榜单以欧美经典剧情片为主，《肖申克的救赎》以 9.7 分、近 200 万评分人数稳居榜首，
体现了"高口碑 + 高参与度"的双重认可，说明评分与评分人数结合能更稳健地衡量影片影响力。

\subsection{Q3：各年代电影数量与平均评分趋势（时间维度）}

\begin{table}[htbp]\centering
\caption{各年代电影数量与平均评分趋势}\label{tab:q3}
\begin{tabular}{lcc}
\toprule
年代 & 电影数量 & 平均评分 \\
\midrule
1950s & 584  & 7.41 \\
1970s & 1174 & 6.97 \\
1990s & 2382 & 7.09 \\
2000s & 6108 & 6.70 \\
2010s & 9885 & 6.08 \\
\bottomrule
\end{tabular}
\end{table}

\textbf{分析}：电影数量随年代快速增长，2010s 收录近万部，是 1990s 的 4 倍多，反映影视
产业的爆发式扩张与数据收录的时效性偏向；但平均评分却从 1950s 的 7.41 持续下滑到 2010s
的 6.08。这一"量升质降"趋势，既可能源于近年商业片、网络电影大量涌入拉低均分，也存在
经典老片经时间筛选后"幸存者偏差"导致老片均分偏高的因素。

\subsection{Q4：每个类型评分最高的电影（窗口函数）}

用 \lstinline!ROW_NUMBER() OVER (PARTITION BY genre ORDER BY rating_score DESC)!
取每个类型评分第一的电影：剧情/犯罪→肖申克的救赎(9.7)、爱情/同性→霸王别姬(9.6)、
悬疑→控方证人(9.6)、喜剧→美丽人生(9.5)、历史/战争→辛德勒的名单(9.5)、动作→这个杀手
不太冷(9.4)、冒险→星际穿越(9.3)。

\textbf{分析}：窗口函数在不缩减数据行的前提下完成"分组内排名"，相比 GROUP BY 更灵活。
结果显示同一部经典片常在多个类型中夺冠（如肖申克的救赎同属剧情与犯罪），反映豆瓣的
多标签体系；各类型头部影片评分都在 9.3 以上，是该类型的标杆之作。

\subsection{Q5：高于本年代平均分的经典电影（JOIN）}

将电影表与"各年代均分"子表 JOIN，筛选评分人数 $\geq 10$ 万且高于本年代均分的电影。
结果如肖申克的救赎(9.7 vs 年代均分 7.09)、千与千寻(9.4 vs 6.70)、忠犬八公的故事
(9.4 vs 6.70) 等。

\textbf{分析}：JOIN 将明细数据与聚合指标关联，实现"相对基准"的衡量。这些影片不仅绝对
分高，更显著高于同年代平均水平（高出 2--2.7 分），是真正"超越时代"的口碑标杆；该方法
比单纯按绝对分排序更能体现影片在其时代背景下的相对卓越程度。

A-2 各查询结果截图见图~\ref{fig:sparka2q123} 与图~\ref{fig:sparka2q45}。

\shot{spark_a2_q123.png}{A-2 查询 Q1--Q3 结果}{fig:sparka2q123}
\shot{spark_a2_q45.png}{A-2 查询 Q4--Q5 结果}{fig:sparka2q45}

\section{A-3：性能对比与 Amdahl 分析}

选取"按年代聚合平均评分"查询，分别用单机 Pandas 与 PySpark（1 个、2 个 Executor）
实现，记录执行时间如表~\ref{tab:perf} 与图~\ref{fig:perf}。

\begin{table}[htbp]\centering
\caption{Pandas vs PySpark 执行时间对比}\label{tab:perf}
\begin{tabular}{lcc}
\toprule
方案 & 执行时间(s) & 相对 Pandas \\
\midrule
Pandas（单机） & 0.42  & 1.0$\times$ \\
PySpark（1 Executor） & 13.58 & 32$\times$ 慢 \\
PySpark（2 Executor） & 15.90 & 38$\times$ 慢 \\
\bottomrule
\end{tabular}
\end{table}

\shot[0.78]{a3_perf_chart.png}{Pandas / PySpark(1) / PySpark(2) 执行时间对比}{fig:perf}

\textbf{Amdahl 定律分析}：本实验有两个反直觉但合理的现象。

\textbf{① 为何 Spark 远慢于 Pandas？} 数据集仅约 39\,MB，属于"小数据"。Pandas 在单机
内存中直接计算，几乎无额外开销；而 Spark 即便处理小数据，也要付出固定的"串行开销"：
SparkContext/Executor 启动、与 OBS 建立 s3a 连接并通过网络读取、任务序列化与分发、
shuffle 落盘与网络传输、结果回收到 Driver 等。按 Amdahl 定律
$S=\frac{1}{(1-p)+p/n}$，当问题规模小时，不可并行的串行部分 $(1-p)$ 占绝对主导，
并行计算节省的那点时间远不足以抵消这些固定开销，故总时间被串行开销"锁死"在十几秒。

\textbf{② 为何 2 个 Executor 反而比 1 个慢？} 加速比仅 0.85$\times$（不升反降）。因为
数据被切分到 2 个 Executor 后，groupBy 聚合需要跨 Executor 进行 shuffle，引入了额外的
网络通信、数据序列化/反序列化与协调开销；而每个 Executor 实际承担的计算量本就微乎其微
（39\,MB/2）。当通信开销 $>$ 并行计算收益时，增加并行度反而拖慢整体——这正是 Amdahl
定律中"通信开销决定并行上限"的体现。

\textbf{结论}：Spark 的价值在于 \textbf{大规模数据}下的水平扩展能力。对本例的小数据集，
单机 Pandas 是更优选择；只有当数据量增长到单机内存无法容纳、串行开销被巨大的并行计算
收益摊薄时，Spark 的分布式优势才会显现。选择计算框架必须匹配数据规模。
